\scanpage{471}
\mspage{238}

pas à la situation, le soumettre peut-être à Carlos, qui est à l'aise \unsure{tant} avec les groupes discrets et leurs générateurs et relations, qu'avec la théorie des \unsure{$\frac{1}{2}$ simples} -- et la \unsure{\uline{topologie}|topologie} ....

\scanpage{472}
\mspage{239}

\uline{33} Étude des revêtements finis -- relation entre les $\unsure{\mathcal N|N}_{g,\nu},\Gamma_{g,\nu}$ pour $g$ variable.

Soit $(\pi\deleted{,\unsure{a}})$ groupe ext. à lacets profini arithmétisé de type $(g,\nu)$, $\pi'$ un \unsure{ss-grpe|ss-grp} ouvert de $\pi$, d'où une application \unsure{régulière}
\[
\unsure{\mathrm{Discr\acute{e}t}}(\pi)\unsure{\hookrightarrow|\longleftarrow}
\unsure{\mathrm{Discr\acute{e}t}}(\pi'),
\]
on voudrait prouver qu'elle est compatible avec la relation d'équivalence \deleted{\unsure{des}\ \illegible{1}}\inserted{\unsure{de l'arithmétisabilité}}, de \deleted{\unsure{prouver}} sorte que \unsure{$\forall e$|tte} arithmétisation de $\pi$ en donne une de $\pi'$, et de $\overline{\mathrm m}$ pour les \deleted{\unsure{discre}} prédiscrétifications.

On voudrait établir en $\overline{\mathrm m}$ temps que les applications
\[
\begin{array}{lll}
 &P^+(\pi)\longrightarrow P^+(\pi')&\begin{gathered}\text{sur les prédiscrétifications}\\\text{orientées}\end{gathered}\\[1em]
\text{et}&A(\pi)\unsure{\longrightarrow}A(\pi')&\text{sur les arithmétisations}
\end{array}
\]
sont injectives. Pour prouver ces points, on est ramené au cas où $\pi'$ est \deleted{\unsure{invariant}} caractéristique. (Cf \S\,28, diagramme 7, \unsure{ff}) \unsure{Choisissant|Choisissons}

\scanpage{473}

une discrétification $\pi_0$ de $\pi$, \unsure{donc} $\pi'_0$ de $\pi'$, on trouve alors
\[
\unsure{\widehat{\widehat{\mathcal G}}}(\pi_0)\longrightarrow
\unsure{\widehat{\widehat{\mathcal G}}}(\pi'_0)
\]
et on a

\textcircled{1} L'hom $\unsure{\widehat{\widehat{\mathcal G}}}(\pi_0)\unsure{\longrightarrow}\unsure{\widehat{\widehat{\mathcal G}}}(\pi'_0)$ envoie $M(\pi_0)$ dans $M(\pi'_0)$

on aura donc un diagramme (\unsure{parent} de celui du \S\,28)
\diagram{[(1) Deux rangées de trois objets. Rangée supérieure, de gauche à droite : $\unsure{\widehat{\mathcal G}^{+}}(\pi_0)$, $M(\pi_0)$, $\unsure{\widehat{\widehat{\mathcal G}}}(\pi_0)$. Rangée inférieure : $\unsure{\widehat{\mathcal G}^{+}}(\pi'_0)$, $M^{\illegible{1}}(\pi'_0)$, $\unsure{\widehat{\widehat{\mathcal G}}}(\pi'_0)$. Chacun des trois objets supérieurs est relié à l'objet inférieur correspondant par une flèche dirigée vers le bas, avec un crochet à son extrémité supérieure. Dans chacune des rangées, les objets voisins sont reliés par un trait horizontal dont l'extrémité gauche est recourbée, ouverte vers la droite ; aucune pointe horizontale n'est nettement discernable. Le signe suivant le $M$ inférieur et le glyphe inférieur droit sont surchargés.]}
qui \deleted{signifie} implique qu'une discrétification $\pi_1$ sur $\pi$ qui donne $\overline{\mathrm m}$ prédiscrétification orientée \unsure{que|par} $\pi_0$ (resp. $\overline{\mathrm m}$ arithmétisation) définit une \inserted{\unsure{discrét}} $\pi'_1$ de $\pi'$ qui donne $\overline{\mathrm m}$ prédiscr. or. (resp. $\overline{\mathrm m}$ arithm.) que $\pi'_0$. Donc on a bien $\unsure{P}^{+}(\pi)\longrightarrow P^+(\pi')$, $A(\pi)\longrightarrow A(\pi')$, il faut encore \unsure{exprimer|examiner} l'injectivité, qui revient aux relations
\[
\begin{gathered}
(2)\quad
\unsure{\widehat{\mathcal G}^{+}}(\pi_0)
=
\unsure{\widehat{\mathcal G}^{+}}(\pi'_0)\cap
\unsure{\widehat{\widehat{\mathcal G}}}(\pi_0),\qquad
M(\pi_0)=M(\pi'_0)\cap
\unsure{\widehat{\widehat{\mathcal G}}}(\pi_0)
\end{gathered}
\]
\scanpage{474}
\mspage{240}

i.e. un \unsure{autom.} à lacets de $\unsure{\pi^{+}|\pi}$ qui, sur $\pi'$, appartient à $\unsure{\widehat{\mathcal G}}(\pi'_0)$ \deleted{\illegible{2}} resp. à $M(\pi'_0)$, est déjà dans $\unsure{\widehat{\mathcal G}}(\pi_0)$ resp. dans $M(\pi_0)$.

L'assertion \deleted{\illegible{1}} \unsure{propre|propré} étant admise, l'assertion \inserted{non} \unsure{propre} signifie aussi que l'hom. can.
\[
(3)\qquad
\unsure{\Gamma}_a=M_a(\pi)/\unsure{\widehat{\mathcal G}}^{+}_a(\pi)
\longrightarrow
\unsure{\Gamma}_{a'}=M_{a'}(\pi')/\unsure{\widehat{\mathcal G}}^{+}_{a'}(\pi')
\]
est \uline{injectif}. Par construction (via $\unsure{\Gamma}_{\unsure{\mathbb Q}}$) c'est aussi surjectif, donc on trouverait :

\textcircled{2} l'hom. canonique (3) est bijectif

et on trouverait aussi une bijection
\[
(4)\qquad P_a(\pi)\longrightarrow P_{a'}(\pi')
\]
entre \unsure{arithm.} de $\pi$ de type $a$, et arithm. de $\pi'$ de type $a'$, \uline{compatible} avec l'iso. \unsure{(3)}.

Enfin, un résultat dans le cas $\pi'$ \unsure{inv.} car. s'étendrait aussitôt au cas d'une $\pi'\subset\pi$ ouvert quelconque.

Procédant \inserted{p.ex.} comme dans le \S 28, à coups de ``correspondances'' \inserted{arithm. extérieures\unsure{,}} entre

\scanpage{475}

``courbes potentielles'', on trouve un groupoïde connexe (ponctué par les $\widehat{\pi}_{g,\nu}$, p.ex, \deleted{\unsure{)}} qui y sont \unsure{can.} isomorphes), dont le $\pi_1$ ext. est la valeur commune des $\unsure{\Gamma}_{g*}=\unsure{\Gamma}_{g,\nu}$.

Pour trouver \unsure{un} l'isom. \inserted{\unsure{canonique}} entre $\unsure{\Gamma}_{g,\nu}$ et $\unsure{\Gamma}_{g',\nu'}$, on choisit une \deleted{\unsure{multiple commun}} \inserted{correspondance entre} un $\pi_{g,\nu}$ et $\pi_{g',\nu'}$, p.ex. en \unsure{regardant} \deleted{\unsure{$N$ de $2g-2+\nu$ et $2g'-2+\nu'$}} l'un et l'autre (quitte à passer : de $g,\nu$ à $g,\widetilde{\nu}$ avec $\widetilde{\nu}\geq\nu$, et de $\overline{\mathrm m}$ pour $g',\nu'$ à $(g',\widetilde{\nu}')$ avec $\widetilde{\nu}'\geq\nu'$) comme des sous-groupes finis du $\overline{\mathrm m}$ $\pi_{0,3}$ \unsure{et} d'où
\[
\unsure{\Gamma}_{0,3}\xrightarrow{\sim}\Gamma_{g,\widetilde{\nu}},
\qquad
\Gamma_{0,3}\xrightarrow{\sim}\Gamma_{g',\widetilde{\nu}'}.
\]

On aimerait maintenant voir ce qui, dans le yoga précédent, est indépendant de notre conjecture. Par exemple, le fait que les \unsure{hom} surjectifs \inserted{\unsure{can.}} $\unsure{\Gamma}_{g,\nu}\unsure{\longrightarrow}\unsure{\Gamma}_{g,\nu'}$ ($g$ le $\overline{\mathrm m}$, $\nu'$ constant) n'est pas établi.

Cependant, si on savait que dans la situation de $(\pi,\pi')$ avec $\pi'$ \unsure{inv.} car. on ait
\[
\unsure{\widehat{\widehat{\mathcal G}}}(\pi'_0)\cap
\unsure{\widehat{\mathcal G}}^{+}(\pi_0)
=
\unsure{\widehat{\mathcal G}}^{+}(\pi_0)
\]

\scanpage{476}
\mspage{241}

(qui est un énoncé de nature relativement anodin sur des groupes à lacets profinis) on conclurait à la \uline{bijectivité} de $\unsure{\Gamma}_{\pi_0}\longrightarrow\unsure{\Gamma}_{\pi'_0}$ dans cette situation, et on en conclurait que le noyau de $\unsure{\Gamma}_{\unsure{\mathbb Q}}\longrightarrow\unsure{\Gamma}_{0,3}$ \deleted{\unsure{s'envoie|l'envoie}} s'envoie \uline{trivialement} dans les $\unsure{\Gamma}_{g,\nu}$ pour $\forall\nu$ -- i.e. qu'il s'envoie dans les $\unsure{\Gamma}_{g,\nu}$ par l'\uline{intermédiaire} de $\unsure{\Gamma}_{0,3}$ \margin{\unsure{$!\! -\! A$}} -- et que de plus $\unsure{\Gamma}_{0,3}\longrightarrow\Gamma_{g,\nu}$ est
\deleted{\unsure{défaut|à défaut} d'avoir établi ceci, \unsure{je vais montrer} au moins que l'image du noyau $\Theta$ de $\unsure{\Gamma}_{\unsure{\mathbb Q}}\longrightarrow\unsure{\Gamma}_{0,3}$ \unsure{ds} les $\unsure{\Gamma}_{g,\nu}$ est \uline{\unsure{finie|fini}} ce qui signifie aussi}
\inserted{\uline{iso} pour $\forall(g,\nu)$} avec $\nu$ \uline{assez} grand -- il suffit que la surface de type $(g,\nu)$ \deleted{\unsure{puisse}} (ou de type $(g,\nu')$ avec un $\nu'\leq\nu$) puisse s'obtenir comme rev. \unsure{étale} de $\unsure{\mathcal U|X}_{0,3}$, i.e. que $\pi_{g,\nu'}$ se réalise \inserted{(\unsure{avec sa str.} à lacets)} comme \unsure{ss-groupe|ss: groupe} \deleted{à lacets} \inserted{d'indice fini} \deleted{\unsure{d'in}} de $\pi_{0,3}$. Alors, \deleted{\unsure{prenant}} \inserted{considérant} un $\pi''$ \inserted{$\unsure{\simeq|\sim}\pi_{g'',\nu''}$} $\unsure{\subset|\in}\pi_{g,\nu'}$ d'indice fini et car. dans $\pi_{0,3}$, on aura
\scanpage{477}

\diagram{\begin{tikzcd}[ampersand replacement=\&]
\unsure{\Gamma}_{0,3} \arrow[r,"\sim"] \& \unsure{\Gamma}_{g'',\nu''} \\
\unsure{\Gamma}_{\mathbb Q} \arrow[u] \arrow[r] \& \unsure{\Gamma}_{g,\nu'} \arrow[u,"\unsure{\wr|2}"']
\end{tikzcd}}
\unsure{diagramme} comm.

d'où notre \unsure{assertion} $\unsure{\Gamma}_{0,3}\xrightarrow{\sim}\unsure{\Gamma}_{g,\nu'}$, et a fortiori l'hom surjectif $\unsure{\Gamma}_{0,3}\longrightarrow\unsure{\Gamma}_{g,\nu}$ qui le \unsure{factorise} \inserted{$(\nu\geq\nu')$} est bijectif. Notons \unsure{p.ex.} que les $\unsure{\Gamma}_{0,\nu}\longrightarrow\unsure{\Gamma}_{0,3}$ $(\nu\geq3)$ sont alors bijectifs. Mais \deleted{il \illegible{3}}\inserted{\deleted{\illegible{1}}} noter qu'en tout état de cause, on ne trouve de résultat pour un $g,\nu$ que si $\nu\geq3$ -- donc ceci ne dit rien sur la fidélité \unsure{éventuelle} \unsure{p.ex.} de $\unsure{\Gamma}_{0,3}\unsure{\longrightarrow|-}\unsure{\Gamma}_{g,0}=\unsure{\Gamma}_g$ $(g\geq2)$ \unsure{ou} $\unsure{\Gamma}_{0,3}\longrightarrow\unsure{\Gamma}_{1,1}=\unsure{\Gamma}_1$.

\margin{On va \unsure{faire} un peu d'\illegible{1} \unsure{conjecturer} \unsure{(établir)} \unsure{pour} l'\unsure{injectivité} du \unsure{(3)}}
Soit une courbe algébrique \inserted{$\unsure{U}$ de type $g,\nu$} définie sur un sous-corps \inserted{$K$} fini sur $\mathbb Q$ de $\overline{\mathbb Q}_0$, \unsure{donc|d'où} elle définit une action extérieure du sous groupe ouvert $\Gamma=\Gamma_K\subset\unsure{\Gamma}_{\mathbb Q}$ sur $\pi_1(\unsure{U}_{\overline{\mathbb Q}_0})$. \deleted{Supposons} Quitte à agrandir $K$, \deleted{\unsure{en} $K'$, \unsure{OPS} \unsure{que}} et à \unsure{agrandir}

\scanpage{478}
\mspage{242}

$\nu$ en $\nu'$ (i.e. faire des trous \inserted{pour \unsure{trouver} $\unsure{U}'$}) \uuline{OPS que} $\unsure{U}'$ \uuline{est un rev. étale de} \unsure{\deleted{$\mathbb P^1_K$}} $(\unsure{U}_{0,3})_K$.
\margin{on pose $\pi'=\pi_1(\unsure{U}'_{\overline{\mathbb Q}_0})$ $\pi=\pi_1(\unsure{U}_{\overline{\mathbb Q}})$}
\deleted{\unsure{Choisissons}} \deleted{\unsure{Considé-rons} $\unsure{\pi|\Pi}$ de $\unsure{U}'$ \unsure{rel}/$K$, donc $\unsure{\Gamma'}$ opère \unsure{effectivement} sur $\unsure{\pi|\Pi}$.}
Soit alors $\unsure{\Theta|\Theta'}$ le noyau de
\diagram{\begin{tikzcd}[ampersand replacement=\&]
\unsure{\Gamma'|\Gamma} \arrow[r] \arrow[d,phantom,"\cap" description] \& \unsure{\Gamma}_{0,3} \\
\unsure{\Gamma}_{\mathbb Q} \&
\end{tikzcd}},
il opère donc \deleted{par} \unsure{autom.|autom} int. sur $\widehat{\pi}_{\unsure{03|0,3}}$, donc \deleted{\unsure{son}}\inserted{l'}image \inserted{de $\Theta'=\Gamma'\unsure{\cap|\wedge}\Theta$} dans le groupe des \unsure{autom.|autom} extérieurs de $\pi'\subset\widehat{\pi}_{0,3}$ est un sous-groupe fini -- a fortiori son image dans $\unsure{\Gamma}_{\pi'}$, et dans \inserted{$\unsure{N}(\pi)$ et dans} $\unsure{\Gamma}_{\pi}$. Il en est de $\overline{\mathrm m}$ (choisissant un point \unsure{base|bon} rat/$K$ dans $\unsure{U}'$) de l'opération \unsure{effective} \unsure{dans|de} $\unsure{\pi_1(U,P)|\pi_1(U',P)}$ -- car il suffit d'appliquer le résultat précédent à $\unsure{U}\unsure{\backslash|-}\{P\}$. On en conclut aussi que \deleted{si $(g,\nu)$ est \unsure{tel}} les noyaux des hom $\unsure{\Gamma}_{0,\nu}\longrightarrow\unsure{\Gamma}_{0,3}$ $(\nu\geq3)$ sont \uline{finis}. On voit facilement que pour \unsure{la|les} V.A. $A$ définie sur $K$, l'opération de $\Theta'$ sur les $\prod_\ell T_\ell(A)$ se fait à travers un groupe quotient fini.

Il \unsure{devient} très difficile \unsure{d'imaginer}, comment il \unsure{\deleted{$\Theta$}} pourrait se faire que $\Theta$ n'opère pas, en fait, trivialement ! \unsure{J'ai}
\scanpage{479}

même l'impression que je peux montrer grâce au résultat du russe que m'a signalé Deligne, \deleted{\unsure{crois je}} que \unsure{$\Gamma_{\mathbb Q}$}\,$\longrightarrow$\,\unsure{$\Gamma_{0,3}$} est iso ! \unsure{Cela|Cel} \unsure{signifierait} que les actions extérieures de \unsure{$\Gamma_{\mathbb Q}$} sur des \unsure{$\pi_1$|$\Pi_1$} \inserted{\unsure{$\simeq\overline{\pi}_{g,\nu}$|$\simeq\overline{\Pi}_{g,\nu}$}} peuvent s'interpréter (au moins pour $\nu$ assez grand, $g$ étant fixé) comme des scindages \unsure{de} l'extension \inserted{\unsure{$N_{g,\nu}$|$\mathcal N_{g,\nu}$}} de \unsure{$\Gamma_{g,\nu}$} par \unsure{$\widehat{\mathcal T}_{g,\nu}$|$\widehat{\mathcal C}_{g,\nu}$} — ou de l'extension \unsure{$M_{g,\nu}$|$\mathcal M_{g,\nu}$} de \unsure{$\Gamma_{g,\nu}$} par \unsure{$\widehat{\mathcal G}_{g,\nu}$|$\widehat{\mathcal C}_{g,\nu}$}, quand il s'agit d'actions effectives.

Bien entendu, la question essentielle qui se pose alors (admettant \unsure{$\Gamma_{\mathbb Q}$}\,$\simeq$\,\unsure{$\Gamma_{0,3}$}) est de \unsure{caractériser} \unsure{$\Gamma_{0,3}$} algébriquement, ainsi que les \unsure{$N_{g,\nu}$|$\mathcal N_{g,\nu}$} — et de donner une description algébrique également des scindages ``admissibles'' — i.e \unsure{correspondant} bel et bien à des courbes algébriques sur des corps de nombres algébriques.

\scanpage{480}
\mspage{243}

\unsure{que|qu'} \deleted{\unsure{il suffit}} \unsure{en} \unsure{tout|tant} \unsure{excluant} les autres relevant l'action ext. \unsure{donnée|donné}, qui ne normalisent aucun \unsure{ss-groupe|ss-groupes} à lacets, \inserted{ou encore telle que \unsure{$\widehat{\pi}_{0,3}^{\Gamma''+}$|$\widehat{\Pi}_{0,3}^{\Gamma''+}$} \deleted{\illegible{1}} $(1)$} (Mais il est \unsure{concevable} qu'il faille y apporter des conditions plus subtiles, faisant intervenir les Frobenius ... ) On peut se proposer de trouver une description des actions \uline{ext.} ``admissibles'', \deleted{\unsure{en termes}} \inserted{sans avoir} à passer par des actions effectives admissibles — et à \unsure{s'endea-|s'ende-}\linebreak\unsure{vorer} d'en donner une définition plus ou moins plausible. La difficulté bien \unsure{sur|sûr} provient du fait que si \unsure{$\Gamma''$} opère \uline{ext.} sur \deleted{\illegible{1}} \unsure{$\widehat{\pi}_{0,3}$|$\widehat{\Pi}_{0,3}$}, il n'opère pas extérieurement pour autant sur un sous-groupe ouvert donné \inserted{\unsure{$\pi'$|$\Pi'$}} de \unsure{$\widehat{\pi}_{0,3}$|$\widehat{\Pi}_{0,3}$}. Mais on va \unsure{supposer} justement qu'il \uline{existe} une telle action, \deleted{\unsure{compatible}} \inserted{de telle façon que} \deleted{\unsure{avec}} \unsure{$\pi'$|$\Pi'$}\,$\longrightarrow$\,\unsure{$\widehat{\pi}_{0,3}$|$\widehat{\Pi}_{0,3}$} soit compatible avec les actions extérieures, et que l'action de \unsure{$\Gamma''$} sur \unsure{$\pi$|$\Pi$} se \unsure{déduise} (mod iso) de
